\section{Многокомпонентные контекстно-свободные языки в качестве ограничений на пути}
	\label{sec:FLPQ_MCFGConstraints}
	\tikzsetfigurename{FLPQ_MCFGConstr_}

\mytodo{Убрать дублирование мотивации с секцией 7}

Создание алгоритмов поиска путей с ограничениями в терминах многокомпонентных контекстно-свободных языков позволяет использовать более сложные графовые запросы. Такие алгоритмы могут найти применение в различных областях, например, в статическом анализе кода. На практике, одним из наиболее широко используемых КС-языков в задачах поиска путей в графе является язык Дика~\sidecite{zhang2013fast,kodumal2004set} правильных скобочных последовательностей. В частности, в статическом анализе программ этот язык используется для \textit{context-sensitive} или \textit{data-dependence} анализа~\sidecite{reps2000undecidability}. А именно,  в \textit{context-sensitive} анализе моделируется сбалансированность вызовов и выходов из процедур с использованием открывающих и закрывающих скобок. Аналогично, в \textit{data-dependence} анализе используется какое-то одно сбалансированное свойство конструкторов языка, например, обращения к полям (т.е. чтения и записи~\sidecite{specificationCFLReachability,10.1145/2001420.2001440}), перенаправления указателей (т.е. ссылки и разыменования~\sidecite{Zheng}) и т.д. Однако точный анализ, охватывающий два или более сбалансированных свойств является неразрешимой задачей~\sidecite{reps2000undecidability}. Например,  \textit{context-sensitive} и \textit{data-dependence} анализ описывается шафлом языков Дика, и этот язык не является контекстно-свободным. Традиционный подход в таком случае заключается в аппроксимации решения с помощью алгоритмов КС-достижимости. Шафл языков Дика можно рассматривать как пересечение двух КС-языков. Однако КС-языки не замкнуты относительно операции пересечения~\sidecite{Hopcroft+Ullman/79/Introduction}. Поэтому приходится жертвовать точностью для одного из сбалансированных свойств, т.е. один из языков Дика для \textit{context-sensitive} или \textit{data-dependence} анализа аппроксимируется регулярным языком~\sidecite{huang2015scalable,sridharan2006refinement}.

Однако для более точного анализа можно использовать и другие классы формальных языков. Например, \textit{линейные конъюнктивные языки} (LCL) могут применяться для \textit{context-sensitive} и \textit{data-dependence} анализа, причём демонстрируют значительные преимущества такого подхода в точности и масштабируемости~\sidecite{10.1145/3093333.3009848}. Таким образом, класс \textit{многокомпонентных контекстно-свободных языков} (MCFL) также может содержать языки, которые могут быть использованы для повышения точности решения некоторых задач статического анализа программ программного. Одним из кандидатов на место такого языка является язык $O_n$, который может моделировать согласованное количество открывающих и закрывающих скобок для \textit{context-sensitive} и \textit{data-dependence} анализа. Эти языки являются аппроксимациями шафлов языков Дика и, как известно, не являются контекстно-свободными. Например, $O_2=\{\omega \in \{a,\overline{a},b,\overline{b}\}^* \mid |\omega|_a=|\omega|_{\overline{a}} \wedge |w|_b=|w|_{\overline{b}}\}$.


Известно, что задача поиска путей с ограничениями в терминах многокомпонентных контекстно-свободных языков разрешима, поскольку MCFL замкнуты относительно пересечения с регулярными языками (т.е. с графами) и информация о достижимости может быть вычислена путем проверки полученного языка на пустоту, что является разрешимой задачей~\sidecite{SEKI1991191}.

На практике хорошим приёмом для получения высокопроизводительных решений задач анализа графов является выражение наиболее критичных вычислений через операции линейной алгебры, например, через матричные операции~\sidecite{doi:10.1137/1.9780898719918}. Существуют эффективные алгоритмы парсинга для MCFL на основе линейной алгебры, использующие умножения булевых матриц~\sidecite{nakanishi1997efficient,cohen2016parsing} и способные лечь в основу новых алгоритмов MCFL-достижимости. Однако алгоритм в работе~\sidecite{cohen2016parsing} может быть применен только для некоторого подкласса многокомпонентных контекстно-свободных грамматик, называемого \textit{несбалансированными}. Поэтому далее в этой главе будет приведён алгоритм поиска путей с ограничениями в терминах многокомпонентных контекстно-свободных языков, основанный на алгоритме из работы~\sidecite{nakanishi1997efficient}.
